\section*{Chapitre \ref{ch:iv2:played-games} --- Jeux joués}
\markboth{\protect\omTocMark{}{Réponses~: Jeux joués}}{\protect\omTocMark{}{Réponses~: Jeux joués}}

\begin{solution}{iq:iv2:played-games:1}
Il faut garder couvert tout ce qui se trouve sous ce paragraphe. Faire glisser la feuille jusqu'à la première ligne d'arrêt, écrire sa cotation, et ne lire la suite qu'ensuite. Chaque manche correspond à une cotation. La lettre en gras placée après un cours acheteur dans les tableaux donne le résultat pour ce cours~: \textbf{S} signifie que l'autre joueur vend un lot à ce cours acheteur, \textbf{B} signifie qu'il achète un lot au cours vendeur, et \textbf{P} signifie qu'il passe. Il faut repérer son propre cours acheteur et ne lire que cette case, car l'ordre des cases est mélangé afin de garder cachée la valeur de l'autre joueur.

\ivround{1}
\ivsays{Donnez-moi un marché de largeur 2 sur la somme.}
\ivstop
% Generated by code/interviews-2/03-played-games/python/gen_lookup.py. Do not edit. French copy by draft/pilot-quant-fr/ivgen_fr.py.
\ivlookup{\ivcell{12.0}{P}\quad \ivcell{16.5}{S}\quad \ivcell{19.0}{S}\quad \ivcell{20.0}{S}\quad \ivcell{14.0}{S}\quad \ivcell{18.0}{S}\quad \ivcell{13.5}{P}\quad \ivcell{10.0}{B}\quad \ivcell{11.5}{B}\quad \ivcell{17.0}{S}\quad \ivcell{15.0}{S}\quad \ivcell{14.5}{S}\quad \ivcell{17.5}{S}\quad \ivcell{15.5}{S}\quad \ivcell{19.5}{S}\quad \ivcell{12.5}{P}\quad \ivcell{16.0}{S}\quad \ivcell{18.5}{S}\quad \ivcell{11.0}{B}\quad \ivcell{13.0}{P}\quad \ivcell{10.5}{B}}

\textbf{Cotation modèle.} Chaque carte vaut en moyenne 5.5, donc la somme vaut 16.5, et la cotation modèle est $15.5$--$17.5$. Un lot est vendu à 15.5.

\ivsays{Vous venez d'en acheter un. Combien vaut maintenant la somme~?}
La valeur de l'autre joueur ne dépend que de sa carte $c$. Elle vaut $c + 2 \cdot \frac{55 - c}{9} = \frac{7c + 110}{9}$, ce qui donne le tableau suivant~:
\begin{center}\small
\begin{tabular}{@{}l*{10}{c}@{}}
Carte $B$ & 1 & 2 & 3 & 4 & 5 & 6 & 7 & 8 & 9 & 10 \\ \hline
Valeur de l'autre joueur & 13.00 & 13.78 & 14.56 & 15.33 & 16.11 & 16.89 & 17.67 & 18.44 & 19.22 & 20.00
\end{tabular}
\end{center}
Ils ont vendu à 15.5, donc leur valeur est inférieure à 15.5, et $B$ est l'une des cartes 1, 2, 3 ou 4. La moyenne de ces quatre valeurs vaut \textbf{14.17}, soit une baisse de plus de 2 après une seule transaction. Il s'agit de la sélection adverse\index{sélection adverse}~: la transaction elle-même est l'information.

\ivround{2}
\ivsays{Cotez-moi de nouveau. Toujours une largeur de 2.}
Aucune carte n'a été retournée.
\ivstop
% Generated by code/interviews-2/03-played-games/python/gen_lookup.py. Do not edit. French copy by draft/pilot-quant-fr/ivgen_fr.py.
\ivlookup{\ivcell{16.0}{S}\quad \ivcell{11.5}{B}\quad \ivcell{12.5}{P}\quad \ivcell{19.0}{S}\quad \ivcell{14.0}{S}\quad \ivcell{10.0}{B}\quad \ivcell{17.0}{S}\quad \ivcell{13.0}{P}\quad \ivcell{10.5}{B}\quad \ivcell{15.5}{S}\quad \ivcell{14.5}{S}\quad \ivcell{15.0}{S}\quad \ivcell{11.0}{B}\quad \ivcell{18.5}{S}\quad \ivcell{17.5}{S}\quad \ivcell{20.0}{S}\quad \ivcell{12.0}{P}\quad \ivcell{13.5}{P}\quad \ivcell{16.5}{S}\quad \ivcell{19.5}{S}\quad \ivcell{18.0}{S}}

\textbf{Cotation modèle.} Le demi-pas le plus proche de 14.17 donne $13$--$15$. L'autre joueur passe. Passer signifie que sa valeur se situe entre le cours acheteur et le cours vendeur affichés, ce qui exclut la carte 4, dont la valeur 15.33 est supérieure à 15. Restent les cartes 1, 2 ou 3, et une fair value de \textbf{13.78}.

\ivround{3}
\ivsays{La carte $A$ est un 9. Votre marché, s'il vous plaît.}
\ivstop
% Generated by code/interviews-2/03-played-games/python/gen_lookup.py. Do not edit. French copy by draft/pilot-quant-fr/ivgen_fr.py.
\ivlookup{\ivcell{15.0}{P}\quad \ivcell{13.5}{B}\quad \ivcell{18.0}{S}\quad \ivcell{20.0}{S}\quad \ivcell{17.5}{S}\quad \ivcell{14.0}{B}\quad \ivcell{15.5}{P}\quad \ivcell{13.0}{B}\quad \ivcell{16.0}{P}\quad \ivcell{19.5}{S}\quad \ivcell{17.0}{S}\quad \ivcell{19.0}{S}\quad \ivcell{16.5}{P}\quad \ivcell{21.0}{S}\quad \ivcell{14.5}{P}\quad \ivcell{18.5}{S}\quad \ivcell{20.5}{S}}

\textbf{Cotation modèle.} Quelqu'un qui ne regarde que le 9 annonce $9 + 2 \cdot \frac{46}{9} = 19.22$. Cela ignore ce que les transactions ont révélé sur $B$. En les utilisant, la somme vaut $9 + c + \frac{46 - c}{8}$ pour $c = 1, 2, 3$, soit 15.625, 16.5 ou 17.375, dont la moyenne est \textbf{16.5}. La cotation modèle est $15.5$--$17.5$, et l'autre joueur passe. Les trois valeurs se situent toutes dans cette cotation, si bien que le fait qu'il passe n'apporte rien de nouveau.

\ivround{4}
\ivsays{La carte $C$ est un 6. Cotez maintenant avec une largeur de 1.}
\ivstop
% Generated by code/interviews-2/03-played-games/python/gen_lookup.py. Do not edit. French copy by draft/pilot-quant-fr/ivgen_fr.py.
\ivlookup{\ivcell{16.5}{P}\quad \ivcell{17.5}{S}\quad \ivcell{16.0}{P}\quad \ivcell{14.0}{B}\quad \ivcell{15.0}{B}\quad \ivcell{19.5}{S}\quad \ivcell{20.0}{S}\quad \ivcell{15.5}{B}\quad \ivcell{19.0}{S}\quad \ivcell{18.0}{S}\quad \ivcell{18.5}{S}\quad \ivcell{14.5}{B}\quad \ivcell{17.0}{P}}

\textbf{Cotation modèle.} La somme vaut maintenant $15 + B$, soit 16, 17 ou 18, avec une fair value de \textbf{17}. La cotation modèle est $16.5$--$17.5$. L'autre joueur passe, et cela tranche~: une somme de 16 l'aurait conduit à vendre à 16.5, et une somme de 18 à acheter à 17.5. On sait maintenant que $B = 2$ avant que la carte ne soit retournée.

\medskip
\textbf{Règlement-livraison.} La carte $B$ est un 2, donc la somme vaut 17. Il faut valoriser ses propres transactions à 17. Le tableau compare deux joueurs.
\begin{center}\small
\begin{tabular}{@{}lll@{}}
 & Tient compte des transactions & Ignore les transactions \\ \hline
Manche 1 & $15.5$--$17.5$~: achète à 15.5 & $15.5$--$17.5$~: achète à 15.5 \\
Manche 2 & $13$--$15$~: passe & $15.5$--$17.5$~: achète à 15.5 \\
Manche 3 & $15.5$--$17.5$~: passe & $18$--$20$~: achète à 18 \\
Manche 4 & $16.5$--$17.5$~: passe & $19.5$--$20.5$~: achète à 19.5 \\ \hline
Résultat & $+1.50$ & $-0.50$
\end{tabular}
\end{center}

\iqwrong{Traiter chaque carte retournée comme la seule information disponible. Le joueur de droite continue d'acheter et ne se demande jamais pourquoi chaque transaction est un achat. Lorsque toutes les transactions vont dans le même sens, la fair value est fausse.}
\iqfollow{Vous avez réalisé 1.50. La première cotation était-elle bonne~?}{Non, le 9 l'a sauvée. Sur l'ensemble des transactions, cette cotation est exécutée avec ce joueur dans 8 cas sur 10, et chaque transaction perd en moyenne 1.33, soit un coût moyen de 1.07. Une largeur de 2 est trop serrée face à quelqu'un qui a vu une carte. Dans un jeu réel, il faut élargir la cotation lorsqu'on soupçonne que l'autre côté en sait davantage.}
\iqfollow{Dans la manche 3, auriez-vous pu trouver $B$ tout de suite~?}{Oui, si une cotation de largeur 1 était permise. À $16$--$17$, la carte 1 amène l'autre joueur à vendre et la carte 3 à acheter, tandis que la carte 2 le fait passer. On perd 0.375 dans les deux premiers cas et rien dans le troisième, si bien que l'information coûte 0.25 en moyenne. Payer un peu pour apprendre est un vrai choix, et il vaut la peine de le dire à voix haute.}
\iqfollow{Une vraie personne négocierait-elle ainsi~?}{Non. Une personne peut ne pas calculer exactement et ne négocier que pour tester le candidat. Il faut mettre à jour dans le même sens, mais moins fortement, et le dire.}
\iqwhy{Coter revient à poser des questions. Le recruteur observe si l'on met à jour immédiatement après une transaction, si les calculs sont rapides et si l'on sait expliquer chaque mouvement. Obtenir le plus gros gain d'une seule transaction compte moins.}
\end{solution}

\begin{solution}{iq:iv2:played-games:2}
Il faut garder couvert tout ce qui se trouve sous cette ligne, et faire glisser la feuille d'une ligne d'arrêt à la fois.

\ivround{1}
\ivsays{Cotez-moi un marché sur le call 8, de largeur 0.30 au maximum.}
\ivstop

\textbf{Cotation modèle.} La somme dépasse 8 de 1, 2, 3 ou 4 dans 4, 3, 2 et 1 des 36 issues, donc le call vaut $\frac{4 + 6 + 6 + 4}{36} = \frac59 \approx 0.56$. La cotation modèle est $0.50$--$0.70$. Le client dit~: «~J'achète 10 à votre cours vendeur.~» Si le cours vendeur affiché était au-dessus de 0.80, le client dit~: «~Trop large,~» et achète 10 à 0.80 à la place.

On a vendu au-dessus de la fair value, donc en moyenne le gain est de $10 \times (0.70 - \frac59) = 1.44$. On porte aussi un risque~: si les dés affichent 6 et 6, il faut payer 40.

\ivround{2}
Le contrat à terme de l'autre dealer cote 6.90 en cours acheteur et 7.10 en cours vendeur. Que faites-vous~?
\ivstop

\textbf{Réponse modèle.} Acheter 3 contrats à terme à 7.10. Le ratio de couverture est la pente de la valeur du call par rapport au dé 1. Une fois que le dé 1 affiche $1, 2, \dots, 6$, le call vaut $0, 0, \frac16, \frac12, 1, \frac53$, et le contrat à terme monte de 1 à chaque pas. La pente des moindres carrés\index{couverture en delta} de ces valeurs sur le dé 1 vaut exactement $\frac13$, donc 10 calls nécessitent $3\tfrac13$ contrats à terme, et on arrondit à 3.

\ivround{3}
Le dé 1 est un 5. Le contrat à terme cote 8.40 en cours acheteur et 8.60 en cours vendeur. Valorisez votre portefeuille et couvrez-le.
\ivstop

\textbf{Réponse modèle.} Le contrat à terme vaut $5 + 3.5 = 8.5$. Le call rapporte $0, 0, 0, 1, 2, 3$ selon le dé 2, donc il vaut 1. Les 10 calls vendus à 0.70 perdent 3.00, et les 3 contrats à terme gagnent $3 \times 1.40 = 4.20$, soit un gain net de 1.20. Le nouveau ratio de couverture est la pente de $0, 0, 0, 1, 2, 3$ sur le dé 2, soit $\frac{22}{35} \approx 0.63$. Il en faut 6, donc on en achète 3 de plus à 8.60.

On remarque ce qui s'est passé. Le marché a monté et il a fallu acheter davantage à un prix plus élevé. C'est ce que signifie être short gamma\index{gamma}.

\ivround{4}
L'autre dealer affiche le call à 1.10--1.25, le put à 0.30--0.40 et le contrat à terme à 8.40--8.60, sur jusqu'à 5 lots chacun. Voyez-vous quelque chose~?
\ivstop

\textbf{Réponse modèle.} La parité put--call\index{parité put--call} impose que le call moins le put soit égal au contrat à terme moins 8, soit 0.50. Vendre le call à 1.10 et acheter le put à 0.40 encaisse 0.70. Acheter ensuite le contrat à terme à 8.60 règle le reste~: pour tout $U$, les trois positions rapportent ensemble $-(U - 8) + U - 8.60 = -0.60$. Chaque ensemble verrouille $+0.10$, donc 5 ensembles donnent $+0.50$. L'opération inverse verrouille une perte de 0.55, donc il n'existe qu'un seul sens.

\ivround{5}
Le dé 2 est un 2, donc $U$ vaut 7. Régler les positions.
\ivstop

\begin{center}\small
\begin{tabular}{@{}lr@{}}
10 calls vendus à 0.70, expirent sans valeur & $+7.00$ \\
3 contrats à terme achetés à 7.10, réglés à 7 & $-0.30$ \\
3 contrats à terme achetés à 8.60, réglés à 7 & $-4.80$ \\
5 calls vendus à 1.10, expirent sans valeur & $+5.50$ \\
5 puts achetés à 0.40, versent 1 & $+3.00$ \\
5 contrats à terme achetés à 8.60, réglés à 7 & $-8.00$ \\ \hline
Total & $+2.40$
\end{tabular}
\end{center}
L'arbitrage a rapporté $+0.50$ et la transaction avec le client, avec ses couvertures, a rapporté $+1.90$.

\textbf{La couverture était-elle le bon choix~?} Sans aucune couverture, on aurait gagné $+7.00$ ici, donc sur ce chemin les couvertures coûtent de l'argent. Comparer les trois choix sur l'ensemble des 36 issues, en tenant compte du coût des fourchettes~:
\begin{center}\small
\begin{tabular}{@{}lrr@{}}
Politique & Gain moyen & Écart-type \\ \hline
Sans couverture & 1.44 & 10.39 \\
3 contrats à terme, jamais modifiés & 1.14 & 6.62 \\
3 contrats à terme, puis ajustés après le dé 1 & 0.84 & 3.77
\end{tabular}
\end{center}
La couverture renonce à un peu de gain moyen en échange d'un risque bien moindre. C'est là tout son intérêt, et un seul chemin ne permet pas d'en juger.

\iqwrong{Couvrir en espérant finir dans la monnaie. Cette chance vaut $\frac{5}{18} \approx 0.28$, ce qui n'est proche de la pente de $\frac13$ que par hasard. La pente est ce qui compense la valeur du call lorsque le dé bouge.}
\iqfollow{Après la manche 3, quels dés finaux vous ont le plus fait mal~?}{Avant l'arbitrage, votre gain sur le dé 2, de 1 à 6, vaut $-4.1, +1.9, +7.9, +3.9, -0.1, -4.1$. Vous gagnez le plus lorsque $U$ finit à 8 et vous perdez aux deux extrémités. C'est la forme d'un straddle vendu.}
\iqwhy{C'est ainsi que raisonne un desk d'options à petite échelle~: coter avec un avantage, couvrir le risque, connaître son gamma et repérer un prix qui casse la parité. Les recruteurs regardent surtout si vous annoncez chaque étape avant de l'appliquer.}
\end{solution}

\begin{solution}{iq:iv2:played-games:3}
Gardez tout ce qui se trouve sous cette ligne couvert, et faites glisser la feuille vers le bas, une ligne d'arrêt à la fois. Dans les tableaux, \textbf{W} après un cours acheteur signifie qu'il remporte le lot, et \textbf{L} qu'il le perd. Ne lisez que la case correspondant à votre propre enchère.

\ivround{1}
\ivsays{Enchères scellées, et le gagnant paie sa propre enchère. Votre dé vaut 8. Votre adversaire enchérit son dé plus 1. Votre enchère~?}
\ivstop
% Generated by code/interviews-2/03-played-games/python/gen_lookup.py. Do not edit. French copy by draft/pilot-quant-fr/ivgen_fr.py.
\ivlookup{\ivcell{5.0}{W}\quad \ivcell{19.0}{W}\quad \ivcell{1.5}{L}\quad \ivcell{17.0}{W}\quad \ivcell{4.0}{L}\quad \ivcell{11.5}{W}\quad \ivcell{0.5}{L}\quad \ivcell{19.5}{W}\quad \ivcell{18.0}{W}\quad \ivcell{15.5}{W}\quad \ivcell{6.5}{W}\quad \ivcell{3.5}{L}\quad \ivcell{14.5}{W}\quad \ivcell{12.0}{W}\quad \ivcell{2.0}{L}\quad \ivcell{2.5}{L}\quad \ivcell{8.5}{W}\quad \ivcell{7.0}{W}\quad \ivcell{16.5}{W}\quad \ivcell{8.0}{W}\quad \ivcell{0.0}{L}\quad \ivcell{13.0}{W}\quad \ivcell{13.5}{W}\quad \ivcell{1.0}{L}\quad \ivcell{3.0}{L}\quad \ivcell{5.5}{W}\quad \ivcell{10.0}{W}\quad \ivcell{9.5}{W}\quad \ivcell{12.5}{W}\quad \ivcell{6.0}{W}\quad \ivcell{16.0}{W}\quad \ivcell{17.5}{W}\quad \ivcell{15.0}{W}\quad \ivcell{10.5}{W}\quad \ivcell{18.5}{W}\quad \ivcell{20.0}{W}\quad \ivcell{7.5}{W}\quad \ivcell{11.0}{W}\quad \ivcell{4.5}{W}\quad \ivcell{14.0}{W}\quad \ivcell{9.0}{W}}

\textbf{Enchère modèle~: 8.5.} Pris isolément, votre 8 rend le lot égal à $8 + 5.5 = 13.5$, et ce nombre est le piège. Vous ne gagnez que lorsque le dé de votre adversaire est bas, donc ce qui compte, c'est la valeur du lot sachant que vous gagnez. Une enchère de 8.5 bat les dés 1 à 7, dont la moyenne est 4, si bien que le lot vaut alors 12 en moyenne, et non 13.5. C'est la malédiction du vainqueur\index{malédiction du vainqueur}. Le gain espéré est la probabilité de gagner multipliée par la marge moyenne lorsque vous gagnez~:
\begin{center}\small
\begin{tabular}{@{}l*{7}{c}@{}}
Enchère & 6.5 & 7.5 & 8.5 & 9.5 & 10.5 & 11.5 & 13.5 \\ \hline
Probabilité de gagner & 0.5 & 0.6 & 0.7 & 0.8 & 0.9 & 1.0 & 1.0 \\
Gain espéré & 2.25 & 2.40 & 2.45 & 2.40 & 2.25 & 2.00 & 0.00
\end{tabular}
\end{center}
Il faut donc enchérir 5 en dessous de votre valeur naïve. C'est la sous-enchère stratégique\index{sous-enchère stratégique}. Le gain reste plat près du meilleur cours acheteur, puis chute vite lorsqu'on remonte vers la valeur naïve. Le dé de votre adversaire valait 3, il enchérit donc 4. Une enchère de 8.5 gagne et paie 8.5 pour un lot qui vaut 11, soit un gain de 2.50.

\ivround{2}
\ivsays{Enchères scellées à nouveau, et le gagnant paie désormais l'enchère du perdant. Votre dé vaut 3. Votre adversaire enchérit le double de son dé. Votre enchère~?}
\ivstop
% Generated by code/interviews-2/03-played-games/python/gen_lookup.py. Do not edit. French copy by draft/pilot-quant-fr/ivgen_fr.py.
\ivlookup{\ivcell{13.0}{W}\quad \ivcell{2.0}{L}\quad \ivcell{4.5}{L}\quad \ivcell{13.5}{W}\quad \ivcell{11.5}{W}\quad \ivcell{3.5}{L}\quad \ivcell{1.0}{L}\quad \ivcell{0.5}{L}\quad \ivcell{15.0}{W}\quad \ivcell{4.0}{L}\quad \ivcell{5.5}{L}\quad \ivcell{2.5}{L}\quad \ivcell{8.0}{L}\quad \ivcell{15.5}{W}\quad \ivcell{11.0}{W}\quad \ivcell{10.5}{W}\quad \ivcell{7.5}{L}\quad \ivcell{9.0}{W}\quad \ivcell{14.0}{W}\quad \ivcell{3.0}{L}\quad \ivcell{17.5}{W}\quad \ivcell{7.0}{L}\quad \ivcell{20.0}{W}\quad \ivcell{8.5}{W}\quad \ivcell{18.0}{W}\quad \ivcell{17.0}{W}\quad \ivcell{6.5}{L}\quad \ivcell{16.0}{W}\quad \ivcell{19.5}{W}\quad \ivcell{9.5}{W}\quad \ivcell{16.5}{W}\quad \ivcell{10.0}{W}\quad \ivcell{5.0}{L}\quad \ivcell{19.0}{W}\quad \ivcell{18.5}{W}\quad \ivcell{14.5}{W}\quad \ivcell{12.0}{W}\quad \ivcell{1.5}{L}\quad \ivcell{0.0}{L}\quad \ivcell{12.5}{W}\quad \ivcell{6.0}{L}}

\textbf{Enchère modèle~: 6.} Ce que vous payez ne dépend plus de votre propre enchère, donc vous ne devez pas sous-enchérir pour cette raison. Votre enchère décide seulement quand vous gagnez. Vous voulez gagner exactement lorsque le prix est inférieur à la valeur. Si le dé de votre adversaire vaut $B$, vous paieriez $2B$ pour un lot qui vaut $3 + B$, ce qui est un gain lorsque $B$ est inférieur à 3. Enchérissez donc le double de votre propre dé, soit 6. En moyenne, vous gagnez $(2 + 1)/10 = 0.30$. Une enchère à votre valeur naïve de 8.5 gagne aussi lorsque $B = 3$ et $B = 4$, où elle rapporte 0 et $-1$, si bien qu'elle rapporte 0.20 en moyenne.

Le dé de votre adversaire valait 4, il enchérit donc 8. Une enchère de 6 perd et vous ne gagnez rien. Une enchère de 8.5 paie 8 pour un lot qui vaut 7, soit une perte de 1. Perdre une enchère qu'il fallait perdre est un bon résultat.

\ivround{3}
\ivsays{Le prix part de 0 et monte. Votre dé vaut 8. Votre adversaire reste en lice jusqu'à deux fois son dé. Quelle est votre limite~?}
\ivstop
% Generated by code/interviews-2/03-played-games/python/gen_lookup.py. Do not edit. French copy by draft/pilot-quant-fr/ivgen_fr.py.
\ivlookup{\ivcell{4.5}{L}\quad \ivcell{19.0}{W}\quad \ivcell{1.0}{L}\quad \ivcell{13.5}{L}\quad \ivcell{2.5}{L}\quad \ivcell{3.0}{L}\quad \ivcell{9.0}{L}\quad \ivcell{11.5}{L}\quad \ivcell{4.0}{L}\quad \ivcell{19.5}{W}\quad \ivcell{12.0}{L}\quad \ivcell{10.0}{L}\quad \ivcell{8.0}{L}\quad \ivcell{5.0}{L}\quad \ivcell{6.5}{L}\quad \ivcell{6.0}{L}\quad \ivcell{9.5}{L}\quad \ivcell{7.0}{L}\quad \ivcell{5.5}{L}\quad \ivcell{12.5}{L}\quad \ivcell{14.0}{L}\quad \ivcell{15.0}{W}\quad \ivcell{13.0}{L}\quad \ivcell{16.0}{W}\quad \ivcell{15.5}{W}\quad \ivcell{0.0}{L}\quad \ivcell{2.0}{L}\quad \ivcell{18.5}{W}\quad \ivcell{10.5}{L}\quad \ivcell{3.5}{L}\quad \ivcell{0.5}{L}\quad \ivcell{17.0}{W}\quad \ivcell{16.5}{W}\quad \ivcell{11.0}{L}\quad \ivcell{20.0}{W}\quad \ivcell{1.5}{L}\quad \ivcell{18.0}{W}\quad \ivcell{8.5}{L}\quad \ivcell{7.5}{L}\quad \ivcell{14.5}{W}\quad \ivcell{17.5}{W}}

\textbf{Limite modèle~: 16.} C'est la même chose que la manche 2, à l'apparence près. Vous payez le prix auquel votre adversaire se retire, et votre limite ne décide que de savoir si vous êtes encore là. Demandez-vous donc, à chaque prix $p$~: si l'adversaire se retire tout de suite, le lot vaut-il plus que $p$~? Se retirer au prix $p$ signifie $B = p/2$, donc le lot vaut $8 + p/2$. Cela dépasse $p$ tant que $p$ est inférieur à 16.

Le dé de votre adversaire valait 7, il se retire donc à 14. Une limite de 16 gagne à 14, et le lot vaut 15, soit un gain de 1. Une limite de 13.5, la valeur naïve, se retire en premier et ne rapporte rien. Ici, l'enchérisseur naïf est trop timide~: les dés élevés demandent une limite au-dessus de la valeur que l'on avait d'abord envisagée, car un adversaire encore en lice vous a signalé que son dé est élevé.

\ivround{4}
\ivsays{Comme dans la manche 1~: votre adversaire enchérit son dé plus 1, et le gagnant paie sa propre enchère. Votre dé vaut 5. Payez-vous 2 pour voir son dé avant d'enchérir~?}
\ivstop
\textbf{Réponse modèle~: payer.} Sans regard, la meilleure enchère est 5.5, qui gagne sur les dés 1 à 4, et rapporte $0.4 \times (5 + 2.5 - 5.5) = 0.80$. Avec le regard, vous enchérissez chaque fois 0.5 au-dessus de son enchère, et vous gagnez toujours un lot qui vaut $5 + B$ pour $B + 1.5$, ce qui rapporte 3.5 quel que soit le dé. Le regard vaut donc $3.5 - 0.8 = 2.70$. Il coûte 2, vous payez donc, et vous espérez $1.50$ au total.

\ivsays{Il vaut 6, il enchérit donc 7. Votre enchère~?}
\ivstop
\textbf{Enchère modèle~: 7.5.} Vous payez 7.5 pour un lot qui vaut 11, vous gagnez donc 3.5, et 1.50 après avoir payé le regard. Remarquez que les 3.5 étaient fixés avant que vous ne voyiez le dé.

\medskip
\textbf{Règlement-livraison.} Le tableau compare le modèle à un joueur qui enchérit ou reste en lice toujours à la valeur naïve du lot, $A + 5.5$, et qui ne paie jamais pour un regard.
\begin{center}\small
\begin{tabular}{@{}lrrrr@{}}
 & \multicolumn{2}{c}{Cette donne} & \multicolumn{2}{c}{Moyenne sur le dé de l'adversaire} \\
 & Modèle & Naïf & Modèle & Naïf \\ \hline
Manche 1 & $+2.50$ & $-2.50$ & $+2.45$ & $0.00$ \\
Manche 2 & $0.00$ & $-1.00$ & $+0.30$ & $+0.20$ \\
Manche 3 & $+1.00$ & $0.00$ & $+2.80$ & $+2.70$ \\
Manche 4 & $+1.50$ & $+0.50$ & $+1.50$ & $-0.45$ \\ \hline
Total & $+5.00$ & $-3.00$ & $+7.05$ & $+2.45$
\end{tabular}
\end{center}
Une seule donne prouve peu de chose~: les dés ont favorisé le modèle de 8. Les colonnes de droite constituent la comparaison équitable, car elles fixent vos propres dés à ceux de cette donne, si bien que la moyenne sur toutes les donnes serait différente elle aussi.

\iqwrong{Enchérir ce que le lot vaut pour vous, sur la seule base de votre information. À la manche 1, cela donne un gain nul. Aux manches 2 et 3, l'enchère est trop haute avec un dé faible et trop basse avec un dé élevé. La question est toujours de savoir ce que vaut le lot sachant que vous gagnez, ou sachant que vous venez de gagner à ce prix.}
\iqfollow{Pourquoi la meilleure enchère dépend-elle autant de qui paie quoi~?}{Prenez des dés uniformes sur 0 à 1. Face à un adversaire qui enchérit son propre dé, votre meilleure enchère au premier prix est votre propre dé. Face à un adversaire qui enchérit le double de son dé, votre meilleure enchère au second prix est le double du vôtre. Payer sa propre enchère la divise par deux.}
\iqfollow{Un véritable recruteur jouerait-il selon une règle fixe~?}{Non. L'adversaire serait une personne, ou plusieurs, et la règle serait quelque chose que l'on estime à partir de ses enchères. On dit alors ce que l'on suppose, et on le met à jour après chaque manche.}
\iqwhy{Les enchères apparaissent partout où un actif est vendu à plusieurs acheteurs à la fois, et remporter une enchère est en soi une information. Le recruteur observe si vous vous demandez ce que la victoire vous apprend avant d'annoncer un montant.}
\end{solution}

\begin{solution}{iq:iv2:played-games:4}
Garder tout ce qui se trouve sous cette ligne couvert, et faire glisser la feuille vers le bas, une ligne d'arrêt à la fois.

\ivround{1}
\ivsays{Une pièce à 60~\% à la mise égale, et votre capital de jeu est de 1\,000. Quelle fraction misez-vous~?}
\ivstop
\textbf{Modèle~: 20~\%, soit 200.} Pour un pari qui rapporte $b$ fois la mise, la fraction de Kelly\index{critère de Kelly} vaut $f = p - (1 - p)/b$, ce qui donne ici $0.6 - 0.4 = 0.2$. La pièce perd. Vous êtes maintenant à 800.

\ivround{2}
\ivsays{Un dé équilibré, et un six rapporte 4 fois votre mise. Quelle fraction~?}
\ivstop
\textbf{Modèle~: rien.} Chaque unité misée vaut $4/6 - 5/6 = -1/6$. La formule donne $f = 1/6 - (5/6)/4 < 0$. Un six sort. Vous restez à 800, et vous aviez raison.

\ivround{3}
\ivsays{45~\% de chances de gagner, avec un paiement de 2 fois la mise. Le pari gagne moins d'une fois sur deux. Quelle fraction~?}
\ivstop
\textbf{Modèle~: 17.5~\%, soit 140.} Le taux de réussite n'est pas l'avantage. Chaque unité vaut $0.45 \times 2 - 0.55 = 0.35$, et $f = 0.45 - 0.55/2 = 0.175$. Le pari gagne, et vous gagnez 280. Vous êtes maintenant à 1\,080.

\ivround{4}
\ivsays{Deux pièces indépendantes à 60~\% à la mise égale, avec la même mise sur chacune. Quelle fraction sur chacune~?}
\ivstop
\textbf{Modèle~: environ 19~\% sur chacune, soit environ 205 sur chacune.} Les traiter une à une donne 20~\%, ce qui est presque juste, puisqu'elles sont indépendantes. La réponse exacte vient de la maximisation de $0.36\ln(1 + 2f) + 0.16\ln(1 - 2f)$, qui donne $f = 0.192$. S'il s'agissait de la même pièce, vous miseriez 10~\% sur chacune, car il n'y aurait alors qu'un seul pari. L'une gagne et l'autre perd, vous restez donc à 1\,080.

\ivround{5}
\ivsays{90~\% de chances de gagner votre mise, et 10~\% de perdre 5 fois votre mise. Quelle fraction~?}
\ivstop
\textbf{Modèle~: 8~\%, soit environ 86.} La valeur par unité vaut $0.9 - 0.5 = 0.4$, ce qui semble excellent. Une mise de 20~\% mène à la ruine. La maximisation de $0.9\ln(1 + f) + 0.1\ln(1 - 5f)$ donne $0.9(1 - 5f) = 0.5(1 + f)$, donc $f = 0.08$. La règle de la moyenne sur la variance donne $0.4/3.24 = 12~\%$, ce qui est trop élevé à cause de la grande perte. Le pari gagne. Vous finissez à 1\,166.40.

\medskip
\textbf{Règlement-livraison.} Le tableau compare le modèle à un joueur imprudent, qui double chaque mise du modèle et mise 20~\% sur le dé, et à un joueur de demi-Kelly, qui mise la moitié des mises du modèle.
\begin{center}\small
\begin{tabular}{@{}lrr@{}}
 & Capital de jeu, cette donne & Croissance logarithmique espérée, par séance \\ \hline
Modèle & 1\,166 & $+0.107$ \\
Demi-Kelly & 1\,100 & $+0.079$ \\
Imprudent & 2\,130 & $-0.141$
\end{tabular}
\end{center}
Une seule donne prouve peu de chose~: le six du joueur imprudent en fait le vainqueur ici, et pourtant il fond à long terme, d'environ 13~\% par séance. Le joueur de demi-Kelly renonce à environ un quart de la croissance pour moitié moins de risque, ce qui est un bon compromis lorsque les cotes ne vous sont pas données d'avance.

\iqwrong{Maximiser le gain espéré en argent. Cela conseille de tout miser aux manches 1, 3 et 4, et 20~\% à la manche 5, la mise maximale sans devoir plus que l'on possède, et une perte à la manche 1 met fin au jeu. L'autre piège classique consiste à miser à la manche 2 parce qu'un six rapporte 4 fois.}
\iqfollow{Pourquoi ne pas toujours utiliser la fraction de Kelly pleine~?}{Ici, les cotes sont données. En pratique, on les estime, et une surestimation de l'avantage à Kelly plein peut faire tomber la croissance à zéro. La plupart des gens en misent donc une fraction.}
\iqfollow{Un véritable recruteur vous donnerait-il les cotes~?}{Rarement. Il vous demanderait de les estimer, ou les modifierait après votre mise, et pourrait ajouter une limite de perte. Dites ce que vous supposez, et comment vous dimensionneriez si cette supposition était fausse.}
\iqwhy{La taille de la mise est ce qui transforme une bonne idée en perte. Le recruteur observe si vous renoncez à un mauvais pari qui gagne par chance, et si vous remarquez que l'avantage de la manche 5 s'accompagne d'un écart.}
\end{solution}

\begin{solution}{iq:iv2:played-games:5}
Garder tout ce qui se trouve sous cette ligne couvert, et faire glisser la feuille vers le bas, une ligne d'arrêt à la fois.

\ivround{1}
\ivsays{Le nombre de nombres premiers inférieurs à 1\,000. Votre cotation~?}
\ivstop
\textbf{Modèle~: entre 150 et 180.} Le théorème des nombres premiers donne $1\,000/\ln 1\,000 = 145$, et cette estimation est trop basse d'environ 15~\% à cette taille. Le nombre est 168, à l'intérieur de votre cotation, et vous obtenez un score de $-0.26$.

\ivround{2}
\ivsays{Le nombre de chiffres de $100!$. Votre cotation~?}
\ivstop
\textbf{Modèle~: entre 155 et 160.} D'après la formule de Stirling, $\log_{10} 100! \approx 100\log_{10}(100/e) + 1.4 = 156.6 + 1.4 = 158$. Le nombre est 158, à l'intérieur, et vous obtenez un score de $-0.05$. On peut coter serré ici, car on dispose d'une formule.

\ivround{3}
\ivsays{Le nombre de mains de cinq cartes tirées d'un jeu de 52 cartes. Votre cotation~?}
\ivstop
\textbf{Modèle~: entre 2.5 millions et 2.7 millions.} $52 \times 51 \times 50 \times 49 \times 48/120$ vaut environ $3.1 \times 10^8/120 = 2.6$ millions. Le nombre est 2\,598\,960, à l'intérieur, et vous obtenez un score de $-0.11$. Si vous l'aviez retenu avec certitude, vous auriez pu coter encore plus serré.

\ivround{4}
\ivsays{La somme des 100 premiers nombres premiers. Votre cotation~?}
\ivstop
\textbf{Modèle~: entre 21\,000 et 25\,000.} Le $n$-ième nombre premier est voisin de $n \ln n$, donc la somme est voisine de $\int_0^{100} x\ln x\,dx \approx 5\,000 \times 4.1 = 20\,500$. Cette formule est trop basse, puisque le 100\up{e} nombre premier vaut 541 et non 460. On majore donc d'un peu plus de 1.1, ce qui donne environ 23\,000. Le nombre est 24\,133, à l'intérieur, et vous obtenez un score de $-0.25$.

\ivround{5}
\ivsays{La somme des chiffres décimaux de $2^{100}$. Votre cotation~?}
\ivstop
\textbf{Modèle~: entre 129 et 150.} $2^{100}$ comporte 31 chiffres, puisque $100 \times 0.301 = 30.1$. Traitons chaque chiffre comme aléatoire entre 0 et 9, de moyenne 4.5 et de variance 8.25. La somme a alors une moyenne de 139.5 et un écart type de 16. Les quartiles sont à 0.67 écart type de la moyenne. Le nombre est 115, en dessous de votre cours acheteur, et vous obtenez un score de $-0.22 - 0.66 = -0.88$. Un manqué sur cinq n'a rien d'étonnant, puisque les quartiles doivent être manqués la moitié du temps.

\medskip
\textbf{Règlement.} Le tableau présente le modèle, un joueur large, dont le cours vendeur vaut toujours le double de son cours acheteur, avec le nombre à l'intérieur, et un joueur serré qui cote à 1~\% de part et d'autre de l'estimation centrale du modèle.
\begin{center}\small
\begin{tabular}{@{}lrrrrrr@{}}
 & 1 & 2 & 3 & 4 & 5 & Total \\ \hline
Modèle & $-0.26$ & $-0.05$ & $-0.11$ & $-0.25$ & $-0.88$ & $-1.55$ \\
Large & $-1$ & $-1$ & $-1$ & $-1$ & $-1$ & $-5.00$ \\
Serré & $-0.08$ & $-0.03$ & $-0.03$ & $-0.25$ & $-1.09$ & $-1.47$
\end{tabular}
\end{center}
Une seule donne prouve peu de chose. Le joueur serré l'emporte ici, parce que les estimations centrales se trouvent être bonnes, et sur des nombres sans formule elles s'effondreraient. La comparaison équitable est le score espéré sous votre propre incertitude. Avec une croyance normale d'écart type $\sigma$ en doublements, vous obtenez un score de $-2.54\sigma$ en cotant les quartiles, $-3.19\sigma$ sans aucune largeur, et $-4.07\sigma$ en cotant 2 écarts types de part et d'autre.

Pourquoi les quartiles\index{quantile}~? Remontez un peu votre cours vendeur. Cela vous coûte ce montant avec certitude, et vous en économise 4 fois autant chaque fois que le nombre se situe au-dessus du cours vendeur. Vous vous arrêtez donc lorsque la probabilité d'être au-dessus vaut $1/4$. La pénalité fixe le quantile~: une pénalité de 20 demanderait les queues à 5~\%.

\iqwrong{Coter la largeur qui vous paraît confortable. Presque tout le monde est trop étroit sur ce qu'il ignore, et trop large sur ce qu'il peut calculer. Il faut se demander comment on l'obtiendrait, puis fixer la largeur à partir de cela.}
\iqfollow{Le recruteur achète immédiatement à votre cours vendeur. Et maintenant~?}{Il connaît le nombre, qui est donc au-dessus de votre cours vendeur, et votre cotation suivante doit se situer entièrement au-dessus de celui-ci. Il faut se demander combien il paierait avant de vendre à nouveau.}
\iqfollow{Un véritable recruteur vous interrogerait-il sur ces nombres~?}{Il vous poserait plus souvent des questions sur le monde réel, comme la population d'une ville ou la consommation d'un vol, et il négocierait avec vous au lieu de vous noter selon une formule. La méthode reste la même~: trouver un chemin, puis fixer la largeur à partir de la distance dont ce chemin peut s'écarter.}
\iqwhy{Coter sur un nombre inconnu teste si l'on sait ce que l'on sait. Un recruteur observe d'où vient la largeur de la cotation.}
\end{solution}

\begin{solution}{iq:iv2:played-games:6}
Garder tout ce qui se trouve sous cette ligne couvert, et faire glisser la feuille vers le bas, une ligne d'arrêt à la fois.

\ivround{1}
\ivsays{Vous êtes à plat. $X$ est coté 7.5 au cours acheteur et 8 au cours vendeur. $Y$ est à 6 au cours acheteur et 6.5 au cours vendeur. Que négociez-vous~?}
\ivstop
\textbf{Modèle~: vendre 1 $X$ à 7.5 et acheter 1 $Y$ à 6.5.} Chaque jambe a un avantage de 0.5, mais prise seule elle ne suffit pas~: $0.5 - 0.1 \times 35/6 = -0.08$. Ensemble, elles forment une position en $Y - X = C - B$, de sorte que le dé commun $A$ s'annule. L'avantage vaut 1, la variance reste $35/6$, et le score est $1 - 0.58 = +0.42$. Deux lots de chaque donnent $2 - 2.33 = -0.33$.

\ivround{2}
\ivsays{Un client vous vend 3 lots de $X$ à 6. Vous devez les prendre. Seul $Y$ est coté maintenant, à 6.75 au cours acheteur et 7.25 au cours vendeur. Que faites-vous~?}
\ivstop
\textbf{Modèle~: vendre 2 $Y$ à 6.75, ce qui laisse 2 $X$ en position longue et 1 $Y$ en position courte.} Le client vous apporte un avantage de 3. Vous êtes désormais long de 2 $X$ et long de 1 $Y$, et la variance vaut $(x^2 + y^2) \times 35/6 + 2xy \times 35/12$, soit 40.8. Le ratio de couverture\index{ratio de couverture} est la covariance divisée par la variance, soit un demi, donc 2 $X$ réclament 1 $Y$ en position courte. Cela donne une variance de 17.5.
\begin{center}\small
\begin{tabular}{@{}l*{4}{r}@{}}
Position en $Y$ après négociation & $+1$ & 0 & $-1$ & $-2$ \\ \hline
Coût de la négociation & 0 & 0.25 & 0.50 & 0.75 \\
Variance & 40.8 & 23.3 & 17.5 & 23.3 \\
Coût plus 0.1 de la variance & 4.08 & 2.58 & 2.25 & 3.08
\end{tabular}
\end{center}
Une position courte de 2 $Y$ ne couvre pas mieux qu'une position courte nulle, et coûte davantage. Couvrir lot pour lot revient à sur-couvrir.

\ivround{3}
\ivsays{$A$ est relancé~: il tombe sur 2. Les deux justes valeurs valent maintenant 5.5. $Y$ est coté 5.4 au cours acheteur et 5.6 au cours vendeur, et $X$ n'est pas coté. Que faites-vous~?}
\ivstop
\textbf{Modèle~: racheter le 1 $Y$ à 5.6.} Le dé qu'ils partageaient a été lancé, de sorte que $X$ et $Y$ sont indépendants. Votre position courte sur $Y$ ne couvre plus rien~: elle ne fait qu'ajouter de la variance. Une fois $A$ connu, chaque lot a une variance de $35/12$. La conserver donne une variance de 14.6, la clôturer donne 11.7, et l'économie de 0.29 l'emporte sur le coût de 0.10.

Les dés valent alors $B = 5$ et $C = 6$, de sorte que $X$ se règle à 7 et $Y$ à 8.

\medskip
\textbf{Règlement.} La trésorerie du modèle est $7.5 - 6.5 - 18 + 13.5 - 5.6 = -9.1$, et 2 lots de $X$ rapportent 14, soit un total de $+4.90$. Le tableau compare un joueur qui prend les lots du client sans rien couvrir, et un joueur qui reste court d'autant de lots de $Y$ qu'il est long de $X$, jusqu'à la fin.
\begin{center}\small
\begin{tabular}{@{}lrrr@{}}
 & Cette donne & Avantage à la juste valeur & Variance lorsque $A$ est relancé \\ \hline
Modèle & $+4.90$ & $+3.40$ & 17.5 \\
Sans couverture & $+3.00$ & $+3.00$ & 52.5 \\
Lot pour lot & $+1.25$ & $+3.25$ & 23.3
\end{tabular}
\end{center}
Une seule donne prouve peu de chose~: le modèle a gagné 1.50 de plus que son avantage, grâce au 2 sur le dé $A$ et à ce qui a suivi. La comparaison équitable est celle des deux colonnes de droite.

\iqwrong{Couvrir un pour un parce que les instruments se ressemblent, ou garder une couverture alors que ce qu'elle couvrait a disparu. Les deux erreurs viennent de la pensée en noms d'instruments et non en dés.}
\iqfollow{Pourquoi le spread était-il meilleur que l'une ou l'autre jambe à la manche 1~?}{Il a deux fois l'avantage pour la même variance, car vendre $X$ et acheter $Y$ annule le dé commun. Le ratio de Sharpe double, et le budget de risque permet de le prendre.}
\iqfollow{Un véritable recruteur vous donnerait-il la covariance~?}{Non. Il montrerait deux prix qui évoluent ensemble et demanderait ce que l'on négocierait contre quoi. On dirait ce que l'on pense qu'ils partagent. Sur les marchés réels, la corrélation d'une couverture varie aussi, en général au mauvais moment.}
\iqwhy{La plupart des desks négocient une chose contre une autre. Un recruteur observe si l'on sait nommer le risque partagé, dimensionner la couverture et la retirer lorsqu'elle cesse de fonctionner.}
\end{solution}

\begin{solution}{iq:iv2:played-games:7}
Garder tout ce qui se trouve sous cette ligne couvert, et faire glisser la feuille vers le bas, une ligne d'arrêt à la fois. Pour chaque manche, le chemin le plus rapide vers la juste valeur passe par la cote~: chaque pile multiplie la cote d'une bonne pièce par $7/3$, et chaque face la divise par $7/3$.

\ivround{1}
\ivsays{Aucun lancer cette manche. Quelle est votre juste valeur, et cotez-vous~?}
\ivstop
\textbf{Modèle~: 50.0, coter entre 47 et 53.} Rien ne peut vous prendre à contre-pied, donc un client vous paie 3 gratuitement.

\ivround{2}
\ivsays{Un lancer est dû. Juste valeur, et cotez-vous~?}
\ivstop
\textbf{Modèle~: 50.0, retirer.} Une pile porte la valeur à 70 et une face à 30. Dans les deux cas, le lecteur rapide obtient 17 de votre part et le client vous paie 3. Une cotation vaut donc $-14$. Le lancer est une pile, et la valeur vaut 70.

\ivround{3}
\ivsays{Un autre lancer est dû. Juste valeur, et cotez-vous~?}
\ivstop
\textbf{Modèle~: 70.0, retirer.} Une pile porte la valeur à 84.5 et une face à 50. La pile a désormais une probabilité de 58~\%. La cotation vaut $3 - 0.58 \times 11.5 - 0.42 \times 17 = -10.8$. Le lancer est une pile.

\ivround{4}
\ivsays{Un autre lancer est dû. Fair value, et cotez-vous~?}
\ivstop
\textbf{Modèle~: 84.5, retirer.} La cote est $(7/3)^2 = 5.44$, donc la valeur vaut 84.5. Une pile porte la valeur à 92.7 et une face à 70. La cotation vaut environ $-4.5$. Le lancer est une pile.

\ivround{5}
\ivsays{Un autre lancer est dû. Fair value, et cotez-vous~?}
\ivstop
\textbf{Modèle~: 92.7, coter entre 89.7 et 95.7.} Désormais la pièce est presque tranchée, et un lancer déplace moins la valeur. Une pile porte la valeur à 96.7, et une face à 84.5.
\begin{center}\small
\begin{tabular}{@{}l*{5}{r}@{}}
Fair value avant un lancer & 50.0 & 70.0 & 84.5 & 92.7 & 96.7 \\ \hline
Taille attendue du mouvement & 20.0 & 16.8 & 10.5 & 5.4 & 2.5 \\
Valeur de la cotation maintenue & $-14.0$ & $-10.8$ & $-4.5$ & $+0.6$ & $+2.7$
\end{tabular}
\end{center}
Le lancer est une face. La valeur baisse à 84.5. Le lecteur rapide nous vend un lot à 89.7, ce qui coûte 5.22, et le client paie 3. On perd 2.22, et la cotation restait pourtant la bonne.

\ivround{6}
\ivsays{Un lancer de plus. Fair value, et cotez-vous~?}
\ivstop
\textbf{Modèle~: 84.5, retirer.} Une pile et une face se compensent, donc on revient à la cote de la manche 4. Le lancer est une pile, et la valeur vaut 92.7. Le jeu s'arrête.

\medskip
\textbf{Règlement-livraison.} Le tableau compare le modèle à un joueur qui cote toujours, et à un joueur qui retire toujours.
\begin{center}\small
\begin{tabular}{@{}lrr@{}}
 & Cette donne & Espérance à chaque décision, cumulée \\ \hline
Modèle & $+0.78$ & $+3.59$ \\
Toujours retirer & $0.00$ & $0.00$ \\
Toujours coter & $-26.14$ & $-30.19$
\end{tabular}
\end{center}
Une seule donne démontre peu de chose~: la cotation risquée du modèle a croisé le lancer malchanceux. C'est la colonne de droite qu'il faut comparer.

\iqwrong{Traiter chaque lancer comme une information de même poids. Le même lancer déplace la valeur de 20 lorsque les cotes sont égales, et d'environ 5 lorsqu'on est presque sûr. L'autre erreur consiste à maintenir une cotation pendant une information, parce que les clients sont là.}
\iqfollow{Et si l'information n'était pas programmée~?}{On ne pourrait pas retirer à temps. Il faudrait alors en tenir compte en permanence, par une cotation plus large, et savoir à quelle fréquence l'information arrive.}
\iqfollow{Un vrai recruteur jouerait-il ainsi~?}{Il annoncerait rarement la probabilité du signal. Il révélerait des données, et observerait si la cotation bouge de la bonne quantité. Un vrai lecteur rapide n'est pas non plus assuré d'obtenir l'exécution, et une vraie cotation peut être annulée à temps certains jours.}
\iqwhy{Rester coté lorsqu'arrive une information est la principale façon, pour un teneur de marché, de perdre de l'argent. Le recruteur observe si l'on se retire aux bons moments, et si l'on revient lorsque l'information est trop petite pour faire mal.}
\end{solution}

\begin{solution}{iq:iv2:played-games:8}
Il faut garder couvert tout ce qui se trouve sous cette ligne, et faire glisser la feuille vers le bas d'une ligne d'arrêt à la fois. Le modèle vient d'un raisonnement à rebours sur les manches, position par position. C'est la bonne réponse à ce jeu, et l'on peut la retrouver à peu près en se demandant ce que coûterait un lot de plus.

\ivround{1}
\ivsays{Vous êtes à plat. Quel est votre centre~?}
\ivstop
\textbf{Modèle~: 100.} On est à plat, donc rien ne pousse à pencher. La valeur du client est 95, donc le client nous vend un lot à 98. On gagne 2, et l'on est long de 1.

\ivround{2}
\ivsays{Vous êtes long de 1. Quel est votre centre~?}
\ivstop
\textbf{Modèle~: 99, soit un cours acheteur à 97 et un cours vendeur à 101.} Un second lot vaut moins que le premier, car il rapproche de la limite et du coût de clôture. On penche donc vers la vente plutôt que vers l'achat. C'est un skew d'inventaire\index{skew d'inventaire}. La valeur du client est 96, donc il nous vend un lot à 97. On gagne 3, et l'on est long de 2, à la limite.

\ivround{3}
\ivsays{Vous êtes long de 2, à la limite, donc votre cours acheteur est retiré. Quel est votre centre~?}
\ivstop
\textbf{Modèle~: 100, donc seulement un cours vendeur à 102.} À la limite, on n'a pas besoin de pencher. La limite a déjà retiré le cours acheteur, donc un centre plus bas ne ferait que céder de l'argent sur le cours vendeur tant qu'il reste du temps. La valeur du client est 94, donc il ne se passe rien, et l'on reste long de 2.

\ivround{4}
\ivsays{Vous êtes toujours long de 2, avec deux manches restantes. Quel est votre centre~?}
\ivstop
\textbf{Modèle~: 99, donc un cours vendeur à 101.} Le temps manque, et deux lots coûteraient 8 à clôturer, donc un cours vendeur moins cher vaut désormais le coup. La valeur du client est 102, donc il achète un lot à 101. On gagne 1, et l'on est long de 1.

\ivround{5}
\ivsays{Vous êtes long de 1, et c'est la dernière manche. Quel est votre centre~?}
\ivstop
\textbf{Modèle~: 98, soit un cours acheteur à 96 et un cours vendeur à 100.} Sortir à 100 ne rapporte rien mais évite les 4 du coût de clôture, et un lot supplémentaire à 96 reste à l'équilibre face à ce qu'il ajoute. La valeur du client est 101, donc il achète à 100. On ne gagne rien, et l'on est à plat.

\medskip
\textbf{Règlement-livraison.} Le tableau donne le meilleur centre à chaque position, issu d'un calcul complet, puis les deux joueurs.
\begin{center}\small
\begin{tabular}{@{}l*{5}{r}@{}}
Position & $-2$ & $-1$ & 0 & $+1$ & $+2$ \\ \hline
Manches 1 à 3 & 100 & 101 & 100 & 99 & 100 \\
Manche 4 & 101 & 101 & 100 & 99 & 99 \\
Manche 5 & 101 & 102 & 100 & 98 & 99
\end{tabular}
\qquad
\begin{tabular}{@{}lrr@{}}
 & Cette donne & Espérance \\ \hline
Modèle & $+6.00$ & $+1.63$ \\
Centre toujours à 100 & $-4.00$ & $+1.13$
\end{tabular}
\end{center}
Sur cette donne, un joueur qui ne penche jamais achète à 98 deux fois, puis reste à la limite avec des cours vendeurs à 102 qui ne sont jamais exécutés, et paie 8 pour clôturer, ce qui fait $-4$. Une seule donne exagère l'écart, qui vaut environ 0.5 en moyenne.

\iqwrong{Pencher davantage à mesure qu'on se rapproche de la limite, y compris à la limite elle-même. Le penchant le plus fort se situe un pas avant la limite. L'autre erreur consiste à ignorer la limite jusqu'à l'atteindre, alors qu'il est trop tard pour agir à moindre coût.}
\iqfollow{Que mesure réellement le penchant~?}{Le coût d'un lot de plus~: la variation de la valeur du reste du jeu lorsque la position augmente d'un lot. Une limite et un coût de clôture rendent cette valeur courbe, et l'on penche d'environ la moitié de l'écart entre l'achat et la vente.}
\iqfollow{Un vrai recruteur vous donnerait-il une valeur connue~?}{Non. La valeur bougerait, et les clients auraient leurs propres avis. Mais une limite est courante dans ces jeux, et le recruteur observe si l'on commence à pencher tôt, puisque la même logique tient.}
\iqwhy{Tout poste de market making a une limite. Le recruteur observe si les cotations changent avec la position, et si l'on reste serein face à la limite.}
\end{solution}

\begin{solution}{iq:iv2:played-games:9}
Il faut garder couvert tout ce qui se trouve sous cette ligne, et faire glisser la feuille vers le bas d'une ligne d'arrêt à la fois.

\ivround{1}
\ivsays{Quelle est la valeur du straddle~? Il est coté 7 au cours acheteur et 7.5 au cours vendeur. Passez-vous un ordre~?}
\ivstop
\textbf{Modèle~: fair value 8. Acheter 2 à 7.5.} L'annonce est toujours plus grande que les petits mouvements, donc leur somme ne modifie jamais en moyenne la distance au strike. Pour $J = 6$, le straddle vaut 6. Pour $J = 10$, il vaut 10. La moyenne est 8.

Le piège est la règle empirique selon laquelle un straddle vaut environ 0.8 écart type. L'écart type vaut ici $\sqrt{2 + 68} = 8.4$, ce qui donne 6.7, donc on vendrait. Cette règle vaut pour des rendements normaux. Ici, le mouvement final est presque toujours grand, si bien que sa taille moyenne est proche de son écart type.

\ivround{2}
\ivsays{L'action est montée de 1, à 101. Quel est votre delta sur les straddles~? Vendez-vous des actions~?}
\ivstop
\textbf{Modèle~: delta 0. Ne rien faire.} Sur le schéma habituel, un straddle gagne du delta lorsque l'action bouge. Ici, non. Avec l'action à 101, le straddle vaut $\tfrac12(|1 + u + J| + |1 + u - J|)$, qui vaut toujours $J$, tant que $J$ est supérieur à $|1 + u|$. Le straddle vaut toujours 8, et un petit mouvement n'a aucun effet sur lui. Vendre de l'action ne ferait qu'ajouter du risque.

\ivround{3}
\ivsays{L'annonce~: hausse de 6. L'action est à 107. Que vaut désormais votre position~? Que faites-vous~?}
\ivstop
\textbf{Modèle~: chaque straddle vaut 7, donc on a perdu 0.50 sur chacun. Vendre 2 actions à 107.} L'action a bougé de 7, ce que l'on aurait pu souhaiter, et l'on a pourtant perdu de l'argent. On a payé 7.5 un straddle qui vaut désormais 7. La baisse de 8 à 7 vient de l'annonce, arrivée à une taille plus petite. Le risque d'annonce a disparu, c'est l'effondrement de la volatilité (vol crush)\index{effondrement de la volatilité}, et ce qui reste vaut 7. Depuis 101, les quatre annonces possibles auraient laissé le straddle valoir 7, 5, 11 ou 9, dont la moyenne est 8~: le straddle ne subit donc aucune baisse attendue à travers l'annonce. L'effondrement est le même dans tous les cas, et l'on ne perd de l'argent sur lui que lorsque le mouvement est plus petit que ce que l'on a payé.

Maintenant, un straddle n'est que l'action moins 100, sans valeur temps~: $|7 + u|$ vaut 6 ou 8. Chacun a un delta de 1, qui ne changera pas. Deux actions vendues à découvert éliminent donc tout le risque restant, et $2 \times 7 = 14$ est verrouillé.

Le dernier mouvement est $-1$, si bien que l'action clôture à 106.

\medskip
\textbf{Règlement.} Vos straddles rapportent $2 \times 6 = 12$, les actions rapportent 2, et vous avez payé 15. Le total est de $-1.00$. Le tableau compare un joueur qui ne vend pas d'actions et un joueur qui a suivi la règle empirique et vendu 2 straddles à 7.
\begin{center}\small
\begin{tabular}{@{}lrr@{}}
 & Cette transaction & Espérance au tour 1 \\ \hline
Modèle & $-1.00$ & $+1.00$ \\
Achète, sans vendre d'actions & $-3.00$ & $+1.00$ \\
Vend selon la règle empirique & $+2.00$ & $-2.00$
\end{tabular}
\end{center}
Une seule transaction prouve peu de chose~: le vendeur s'est trompé et a gagné de l'argent, parce que la taille est tombée petite. Si la même transaction avait donné $J = 10$, le modèle aurait gagné $+7$ et le vendeur $-6$.

\iqwrong{Penser qu'un straddle rapporte chaque fois que l'action bouge beaucoup. Il rapporte lorsque l'action bouge plus que ce qui a été payé. Une autre erreur consiste à couvrir par habitude le petit mouvement du tour~2.}
\iqfollow{Que vous indique le prix de marché de 7.25~?}{Que le marché anticipe un mouvement d'environ 7.25 à l'annonce, contre vos 8. Avec de vraies options, on comparerait le prix avant et après l'échéance qui encadre l'événement, et l'on trouverait la part de la variance due à l'événement.}
\iqfollow{Un vrai recruteur vous donnerait-il le saut~?}{Il vous donnerait un historique passé des mouvements de tels jours, ainsi que les prix des options, et demanderait ce qui est implicite et si vous le conserveriez. Les vrais mouvements ne sont pas non plus à deux points, si bien que le delta ne sera pas exactement nul au tour~2.}
\iqwhy{Les annonces sont le domaine où les prix d'options s'écartent le plus des règles empiriques. Un recruteur observe si vous travaillez à partir de la distribution, ou à partir d'une formule qui suppose une autre distribution.}
\end{solution}

\begin{solution}{iq:iv2:played-games:10}
Garder tout ce qui se trouve sous cette ligne couvert, et faire glisser la feuille vers le bas, une ligne d'arrêt à la fois.

\ivround{1}
Personne n'a encore coté. Votre dé vaut 7. Que vaut la somme~?
\ivstop
\textbf{Modèle~: 23.5.} C'est $7 + 3 \times 5.5$.

\ivround{2}
$A$ cote 19 à l'achat, 20 à la vente. $B$ cote 21 à l'achat, 22 à la vente. Que faites-vous~?
\ivstop
\textbf{Modèle~: acheter 1 à $A$ à 20 et vendre 1 à $B$ à 21.} Une cotation centrée sur 19.5 signifie que le dé de $A$ vaut 3, et une cotation centrée sur 21.5 signifie que celui de $B$ vaut 5. La valeur équitable est donc désormais $7 + 3 + 5 + 5.5 = 20.5$. Ces deux cotations se croisent aussi~: la vente de $A$ est inférieure à l'achat de $B$. Vous gagnez 1 quoi qu'il arrive, et vous êtes à plat.

\ivround{3}
$C$ cote 22 à l'achat, 23 à la vente. Quelle est votre valeur équitable, et que faites-vous~?
\ivstop
\textbf{Modèle~: valeur équitable 21. Vendre 1 à $C$ à 22.} Un centre de 22.5 signifie que le dé de $C$ plus son bruit vaut 6, donc le dé vaut 3, 6 ou 9, avec des probabilités égales, soit une moyenne de 6. La somme vaut donc $7 + 3 + 5 + 6 = 21$. Vendre à 22 offre un avantage de 1. Une cotation bruitée n'est pas sans valeur. Elle a réduit le dé de $C$ à trois valeurs possibles, et elle a déplacé votre valeur équitable de 0.5.

\ivround{4}
$C$ cote désormais 16 à l'achat, 17 à la vente. Quelle est votre valeur équitable, et que faites-vous~?
\ivstop
\textbf{Modèle~: valeur équitable 18. Acheter 1 à $C$ à 17.} Un centre de 16.5 signifie que son dé plus le nouveau bruit vaut 0. Un dé d'au moins 1 et un bruit d'au moins $-3$ font que le dé ne peut valoir que 3, avec un bruit de $-3$, et cela correspond au tour précédent. La somme vaut donc exactement $7 + 3 + 5 + 3 = 18$. Acheter à 17 rapporte 1 à coup sûr, et clôt aussi votre vente à découvert à 22, si bien que vous êtes à plat. Les deux cotations réunies valent bien plus que chacune isolément.

\medskip
\textbf{Règlement.} La somme vaut 18. Le modèle a gagné $21 - 20 = 1$ sur le croisement, et $22 - 17 = 5$ sur $C$, soit un total de $+6$. L'autre joueur ne lit aucune cotation, valorise la somme à 23.5 pendant toute la partie, et achète à chaque vente inférieure à ce montant.
\begin{center}\small
\begin{tabular}{@{}lrr@{}}
 & Cette transaction & Avantage à l'échange, sur ce qui pouvait être connu \\ \hline
Modèle & $+6.00$ & $+3.00$ \\
Dé propre seulement & $-10.00$ & $-2.00$
\end{tabular}
\end{center}
Une seule transaction prouve peu de chose, même si le petit dé de $C$ a ajouté à chacun des deux écarts. Une fois le tour~4 passé, pourtant, les deux cotations ne laissaient qu'un dé valant 3, si bien que rien dans la dernière transaction ne relevait de la chance.

\iqwrong{Trader sur son propre dé comme si les cotations n'étaient que des prix. Une cotation émise par quelqu'un qui sait quelque chose est une partie de son information, et c'est généralement l'élément le plus précieux sur la table.}
\iqfollow{Et votre propre cotation~?}{Elle révèle votre dé de la même façon. Face à des joueurs qui vous lisent, vous pouvez préférer coter tard, ou trader avec une taille qui ne montre pas ce que vous savez, et c'est un coût réel que d'être le premier.}
\iqfollow{Un vrai recruteur vous ferait-il jouer selon des règles fixes~?}{Non. Les autres joueurs actualisent leurs estimations et peuvent bluffer, si bien que vous estimeriez la règle de chacun à partir de la façon dont ses cotations se comportent. Dites ce que vous supposez à leur sujet, et changez-le lorsqu'ils vous surprennent.}
\iqwhy{Les prix portent l'information d'autrui. Un recruteur observe si vous lisez les cotations pour ce qu'elles impliquent, et si vous repérez un marché croisé lorsqu'il existe.}
\end{solution}

\begin{solution}{iq:iv2:played-games:11}
Garder tout ce qui se trouve sous cette ligne couvert, et faire glisser la feuille vers le bas, une ligne d'arrêt à la fois. À chaque tour, le coût espéré d'un achat immédiat de $k$ lots, suivi d'une répartition égale du reste plus tard, est calculé ci-dessous, mesuré à partir du prix courant.

\ivround{1}
Le prix est 100, et la prévision est $+1.5$. Combien en achetez-vous~?
\ivstop
\textbf{Modèle~: 3.} Les lots achetés plus tard coûtent $f = 1.5$ de plus chacun, en moyenne. Les lots restants sont ensuite répartis aussi également que possible.
\begin{center}\small
\begin{tabular}{@{}l*{3}{r}@{}}
Lots maintenant & 2 & 3 & 4 \\
Impact immédiat & 2.0 & 4.5 & 8.0 \\
Coût prévu de l'attente & 6.0 & 4.5 & 3.0 \\
Impact différé & 3.0 & 1.5 & 1.0 \\ \hline
Total & 11.0 & 10.5 & 12.0
\end{tabular}
\end{center}
Vous payez $300 + 4.5 = 304.5$. Puis le prix monte de $1.5$ plus 1, pour atteindre 102.5.

\ivround{2}
Le prix est 102.5, et la prévision est 0. Il a beaucoup monté depuis que vous avez commencé. Combien en achetez-vous~?
\ivstop
\textbf{Modèle~: 1.} Le prix que vous avez manqué est perdu. Sans prévision, l'impact seul conduit à répartir le reste également~: 1 à chacun des trois tours. Vous payez 103. Le prix baisse de 1, pour atteindre 101.5.

\ivround{3}
Le prix est 101.5, et la prévision est $-2$. Il vous reste 2 lots et deux tours. Combien en achetez-vous~?
\ivstop
\textbf{Modèle~: 0.} Attendre économise 2 par lot. Acheter 1 maintenant et 1 plus tard coûte $0.5 + 0.5 - 2 = -1$, et n'acheter rien maintenant et 2 plus tard coûte $2 - 4 = -2$. Doubler plus tard coûte moins que ce que la prévision fait économiser. Le prix baisse de 2 puis monte de 1, pour atteindre 100.5.

\ivround{4}
Le prix est 100.5, et vous devez en acheter 2.
\ivstop
\textbf{Modèle~: 2, comme il se doit.} Vous payez $201 + 2 = 203$.

\medskip
\textbf{Règlement.} Le modèle a payé $304.5 + 103 + 203 = 610.5$. Le tableau compare une répartition égale de 2, 2, 1, 1, et l'achat des 6 lots d'un coup, ainsi que l'attente pour les acheter tous au tour~4.
\begin{center}\small
\begin{tabular}{@{}lrr@{}}
 & Coût de cette transaction & Espérance, avec ces prévisions \\ \hline
Modèle & $610.5$ & $607.5$ \\
Répartition égale & $612.0$ & $609.0$ \\
Tout au tour 1 & $618.0$ & $618.0$ \\
Tout au tour 4 & $621.0$ & $615.0$
\end{tabular}
\end{center}
Une seule transaction prouve peu de chose, même si ici les deux s'accordent. Le bruit touche tous les lots, et seuls ceux qui restent à acheter y sont exposés, si bien que celui qui répartit également n'est pas trop pénalisé. L'écart vient de l'usage de la prévision.

\iqwrong{Traiter le niveau du prix comme un signal. Un prix qui a monté n'est pas cher, et un prix qui a baissé n'est pas bon marché~: seule la prévision indique ce qui vient ensuite. L'autre erreur consiste à acheter tout de suite par anxiété, ce qui coûte 18 d'impact.}
\iqfollow{Et si vous n'aimez pas le risque~?}{Vous achèteriez un peu plus tôt, car les lots restant à acheter portent le risque de prix. Avec une aversion suffisamment forte, vous achèteriez bien davantage au tour~1. Dites ce que vous arbitrez~: un impact certain contre un prix incertain.}
\iqfollow{Un vrai recruteur vous donnerait-il ces règles~?}{Il donnerait rarement une formule d'impact, et en réalité une part de votre impact reste dans le prix et signale aux autres que vous achetez. Vous diriez donc ce que vous supposez, et comment la réponse change si cette part demeure.}
\iqwhy{Exécuter un ordre est une décision quotidienne dans de nombreux desks. Un recruteur observe si vous découpez l'ordre et réagissez à une prévision, et si vous évitez l'envie de courir après le prix ou d'attendre un repli.}
\end{solution}

\begin{solution}{iq:iv2:played-games:12}
Gardez tout ce qui suit cette ligne couvert, et faites glisser la feuille vers le bas d'une ligne d'arrêt à la fois. Le modèle vient d'un raisonnement à rebours. Vous ne cherchez que la probabilité de terminer devant, rien d'autre. Dites-le d'abord~: un recruteur veut savoir ce que vous cherchez à maximiser avant de miser.

\ivround{1}
\ivsays{Vous êtes en retard de 3. Une pièce équilibrée à cote 1 contre 1. Quelle est votre mise~?}
\ivstop
\textbf{Modèle~: 4.} Si vous gagnez, vous êtes en avance de 1, et vous ne misez plus rien ensuite. Si vous perdez, vous pouvez encore revenir.
\begin{center}\small
\begin{tabular}{@{}l*{5}{r}@{}}
Mise & 0 & 1 & 2 & 3 & 4 \\ \hline
Probabilité de terminer devant & 0.658 & 0.607 & 0.629 & 0.611 & 0.665
\end{tabular}
\end{center}
C'est presque une égalité avec une mise nulle, et les deux battent les petites mises. Les petites mises sont ici les pires~: elles consomment des manches sans combler l'écart. La pièce perd, et vous êtes en retard de 7.

\ivround{2}
\ivsays{Vous êtes en retard de 7. 55~\% à cote 1 contre 1. Quelle est votre mise~?}
\ivstop
\textbf{Modèle~: 4.} Il vous faut gagner 8 points avec des mises d'au plus 4, donc deux gains pleins, ou gagner les trois manches qui restent. Misez 4 maintenant~: un gain vous laisse en retard de 3, avec la manche à 60~\% à venir~: $0.55 \times 0.60 = 0.33$. Ne pas miser laisse les manches 3 et 4 à pleines mises~: $0.40 \times 0.60 = 0.24$. Une mise de 1, 2 ou 3 ne donne que 0.132, car il faut alors gagner les trois. Elle gagne, et vous êtes en retard de 3.

\ivround{3}
\ivsays{Vous êtes en retard de 3. 40~\% à cote 1 contre 1, et une manche à venir à 60~\%. Quelle est votre mise~?}
\ivstop
\textbf{Modèle~: 0.} Si vous ne misez rien, une mise de 4 à 60~\% à la manche suivante la gagne~: une probabilité de 0.60. Une mise de 4 ici donne 0.40, car un gain vous met en avance sans rien avoir à faire, et une perte met fin à la partie. Être en retard ne signifie pas prendre un pari quelconque~: vous prenez le meilleur qui reste, et c'est le suivant. La pièce perd, et vous êtes toujours en retard de 3.

\ivround{4}
\ivsays{Vous êtes en retard de 3. 60~\% à cote 1 contre 1, et c'est la dernière manche. Quelle est votre mise~?}
\ivstop
\textbf{Modèle~: 4.} Une mise de 3 égalise, et une égalité est perdante. Toute mise inférieure à 4 revient à ne rien miser. Elle gagne, et vous terminez en avance de 1.

\medskip
\textbf{Règlement.} Le tableau compare le modèle à un joueur qui maximise les points espérés, misant 4 à 55~\% et 60~\% et rien d'autre, et à un joueur qui mise 1 à chaque manche.
\begin{center}\small
\begin{tabular}{@{}lrrr@{}}
 & Écart, cette donne & Probabilité de terminer devant & Écart final espéré \\ \hline
Modèle & $+1$ & 0.665 & $-2.58$ \\
Maximum de points & $+5$ & 0.330 & $-1.80$ \\
Mise 1 à chaque manche & $-3$ & 0.066 & $-2.90$
\end{tabular}
\end{center}
Une seule donne prouve peu de chose~: les deux premiers gagnent ici, et celui qui joue pour les points a le meilleur score. Sur un grand nombre de donnes, le modèle termine devant deux fois plus souvent, et il paie cela par un écart espéré moins bon~: il sacrifie des points pour gagner le prix.

\iqwrong{Jouer pour les points alors que le prix récompense le rang, ou miser le maximum à chaque manche parce qu'on est en retard. À la manche 3, la mise juste selon «~il faut bien tenter le coup~» est fausse~: de meilleures cotes sont encore à venir.}
\iqfollow{Et si vous étiez en avance de 1~?}{Ne misez rien, même à 60~\%. Vous avez déjà gagné, et un bon pari ne peut que vous faire perdre le prix. C'est l'autre moitié de la leçon.}
\iqfollow{Trading-vous ainsi sur un desk~?}{Non. Un bonus peut ressembler à ce prix, et c'est précisément pour cela qu'une firme le plafonne. La firme recherche un profit espéré pour son risque, donc dites que vous savez faire la différence. Et un vrai recruteur peut observer un jeu de ce type pour voir si vous prenez des risques absurdes~: demandez ce que signifie le score avant de commencer.}
\iqwhy{Les gains façonnent les comportements. Un recruteur observe si vous demandez sur quoi vous êtes évalué, et si vous savez raisonner sur le risque lorsque le gain n'est pas linéaire.}
\end{solution}
